% !TEX root = MultipleGPU.tex
% !TEX program = pdflatex
\documentclass[conference]{IEEEtran}
\IEEEoverridecommandlockouts

\usepackage{amsmath,amssymb,amsfonts}
\usepackage{graphicx}
\usepackage{xcolor}
\usepackage{listings}
\usepackage{booktabs}
\usepackage{float}
\usepackage{url}
\usepackage{hyperref}
\usepackage{titlesec}

\titleformat{\section}{\normalfont\Large\centering}{\thesection}{1em}{}

% CUDA / C++ listing style (matches main report)
\lstdefinestyle{cuda}{
  language=C++,
  basicstyle=\ttfamily\scriptsize,
  keywordstyle=\color{blue}\bfseries,
  commentstyle=\color{green!50!black}\itshape,
  stringstyle=\color{red!70!black},
  numbers=left,
  numberstyle=\tiny\color{gray},
  numbersep=5pt,
  frame=single,
  breaklines=true,
  captionpos=b,
  morekeywords={__global__,__device__,__shared__,__constant__,
                __syncthreads,__syncwarp,atomicAdd,__ldg,
                __launch_bounds__,blockIdx,blockDim,threadIdx,
                gridDim,float3,int3,uint32_t,cudaStream_t,
                cudaSetDevice,cudaDeviceEnablePeerAccess,
                cudaMemcpyPeerAsync,cudaMemcpyAsync,
                cudaEventRecord,cudaEventSynchronize,
                cudaStreamWaitEvent},
  tabsize=2
}
\lstset{style=cuda}

% ============================================================
\begin{document}

\title{GPU-Accelerated Discrete Element Method (DEM) Simulation of a Snowball Avalanche Dynamics on NVIDIA Maxwell\\[0.5em]
\LARGE Attachment 1: Multi-GPU Extension Strategy}

\author{
\IEEEauthorblockN{Maurizio Grisotto}
\IEEEauthorblockA{\textit{Politecnico di Torino}\\
Turin, Italy 2026\\
s336145@studenti.polito.it}
}

\maketitle

\begin{abstract}
This attachment presents a theoretical analysis of extending the Angry Santa
DEM snowball simulation to multiple GPUs. The analysis covers pipeline communication 
requirements, 1D domain decomposition along the slope axis, halo exchange sizing, 
a per-frame multi-stream pipeline, comparison with index-based partitioning, 
concrete code change maps, and an optional CUDA-aware MPI path for multi-node clusters.
\end{abstract}

\begin{IEEEkeywords}
multi-GPU, domain decomposition, halo exchange, CUDA peer-to-peer,
\texttt{cudaMemcpyPeerAsync}, MPI, DEM, particle simulation
\end{IEEEkeywords}

% ============================================================
\section{Pipeline Analysis}
\label{sec:pipeline}

The current single-GPU pipeline executes five main stages per frame across
two CUDA streams. 

\begin{table*}[t]
  \caption{Inter-GPU communication per pipeline stage}
  \label{tab:comm}
  \centering
  \begin{tabular}{clp{10cm}}
    \toprule
    \# & Kernel / Stage & Communication \\
    \midrule
    1 & \texttt{captureFromSnowpackKernel}
      & \textbf{Partial.} Snowpack is read-only; replicated on each GPU
        ($\sim$2.7\,MB for 100\,k particles: 6 float + 1 int = 28\,B).
        Captured mass must be \emph{reduced} across all GPUs before
        \texttt{updateCoreMassAfterCapture}. \\[2pt]
    2 & \texttt{forceTetherKernel}
      & \textbf{None.} Per-particle: gravity, damping, tether. \\[2pt]
    3 & Grid build (hash + sort)
      & \textbf{None.} Each GPU builds its own local grid. \\[2pt]
    4 & \texttt{neighborCollisionKernel}
      & \textbf{Yes.} Particles near subdomain boundaries require
        \emph{halo} data from the adjacent GPU
        (Section~\ref{sec:halo}). \\[2pt]
    5 & \texttt{integrateGroundKernel}
      & \textbf{None.} Per-particle semi-implicit Euler integration. \\
    \bottomrule
  \end{tabular}
\end{table*}

Inter-GPU communication is required at exactly two points per frame: the
halo exchange before \texttt{neighborCollisionKernel}, and a small
captured-mass reduce after \texttt{captureFromSnowpackKernel}. Both are
cheap and maskable behind compute. \\
Each GPU will operate with 3 CUDA streams:

\begin{itemize}
  \item \texttt{simStream}: main physics pipeline (existing)
  \item \texttt{gridStream}: grid construction (existing, concurrent with
        \texttt{forceTetherKernel})
  \item \texttt{commStream} (\emph{new CUDA stream}): pack halo $\to$
        \texttt{cudaMemcpyPeerAsync} $\to$ unpack halo
\end{itemize}

The join before \texttt{neighborCollisionKernel} requires \texttt{simStream}
to wait on both \texttt{evGridDone} (existing) and \texttt{evHaloDone} (new).

% ============================================================
\section{1D Domain Decomposition}
\label{sec:decomp}

\subsection{Geometry}

The ball rolls along the slope in the \textbf{s} direction (downhill). Shell
particles cluster around the ball with radius $r_\text{ball} + r_\text{tether}$;
the spatial distribution is essentially 1D along~$s$. The proposed
partitioning divides the $s$ axis into $G$ slices (one per GPU) with
boundaries $s_0 < s_1 < \cdots < s_G$. GPU~$g$ owns all particles with
$s_g \le \text{posS} < s_{g+1}$. The halo region for GPU~$g$ consists of particles in the adjacent
slices $[s_{g-1} - h, s_g)$ and $[s_{g+1}, s_{g+1} + h)$, where $h$ is the halo width (Section~\ref{sec:halo}).

\subsection{Halo Exchange}
\label{sec:halo}

The halo width must cover the maximum interaction radius of
\texttt{neighborCollisionKernel}. The effective cutoff is $\max(\texttt{cohesionRadius},\,2\cdot
\texttt{particleRadius})$, so the required halo width is:
\[
  h = \max\!\bigl(\texttt{cellSize},\;\texttt{cohesionRadius}\bigr)
\]

The halo width must be computed at runtime from the actual
parameter values, not hardcoded to \texttt{cellSize}.

\subsection{Snowpack Replication and Ball Broadcast}

The snowpack is replicated on each GPU; the only post-capture 
communication is reducing \texttt{h\_captureCount} and \texttt{h\_captureMass} across GPUs.
The \texttt{SnowballState} struct ($\approx 60$\,B) is broadcast to all
GPUs via \texttt{cudaMemcpyAsync} (host$\to$device) at the start of each
frame at negligible cost.

% ============================================================
\section{Comparison with Index-Based Partitioning}
\label{sec:index}

The simpler alternative is to assign to each GPU~$g$ the particles with index
$i \bmod G = g$ (round-robin), or contiguous blocks
$[gN/G, (g+1)N/G)$.
With index-based partitioning, particles on the same GPU are not spatially
localised: a particle's neighbors in the hash grid can reside on any other
GPU. This has two severe consequences:
\begin{itemize}
  \item \texttt{O(N)} communication: each GPU must send/receive all positions and velocities
        to/from every other GPU, because it cannot know in advance which particles
        will be neighbors in the spatial grid.
  \item Duplicated grid: each GPU must build a grid over all particles (not only its own),
        because the neighbor search requires global positions. This eliminates the
        advantage of partitioning the grid build.
\end{itemize}

Index-based partitioning requires each GPU to exchange \emph{all} positions
with every other GPU and to build a grid over \emph{all} particles. Domain 
decomposition exploits the strong spatial coherence of shell particles around the ball.

% ============================================================
\section{Code Change Map}
\label{sec:code}

\subsection{New Data Structure}

\begin{lstlisting}[caption={Per-GPU state struct}]
struct PerGPUData {
    int            deviceId;
    ParticleSystem shell;     // N/G particles + halo space
    GridData       grid;      // sized for local + halo
    Snowpack       snowpack;  // replicated read-only copy
    cudaStream_t   simStream, gridStream;
    cudaStream_t   commStream;  // NEW
    cudaEvent_t    evFork, evGridDone;
    cudaEvent_t    evHaloDone;  // NEW
    // Halo buffers: pos(xyz) + vel(xyz) + mass + radius + state
    char  *haloSendBuf, *haloRecvBuf;
    int    haloSendCount, haloRecvCount;
    // Capture accumulators (pinned host + device)
    int   *d_captureCount; float *d_captureMass;
    int   *h_captureCount; float *h_captureMass;
};
\end{lstlisting}

\subsection{New Kernels}

\begin{itemize}
  \item \texttt{packHaloKernel}: scans particles with
        $|\text{posS} - s_\text{border}| < h$ and packs them into a
        contiguous buffer (8 float + 1 int per particle).
  \item \texttt{unpackHaloKernel}: appends received halo particles at
        indices $[\texttt{shellN\_local},\,\texttt{shellN\_local} +
        \texttt{haloRecvCount})$ of the local \texttt{ParticleSystem}, so
        \texttt{neighborCollisionKernel} finds them in the spatial grid.
\end{itemize}

\subsection{Grid Sizing for Halo Particles}

\begin{lstlisting}
// allocateGrid signature: (GridData&, int capacity, const SimParams&)
int maxHaloSize = (int)(shellLocalCap * 0.10f); // 10% headroom
allocateGrid(grid, shellLocalCap + maxHaloSize, params);

// Halo buffer size per border:
size_t haloBytes = maxHaloParticles *
    (8 * sizeof(float) + sizeof(int));
// 8 floats: posX/Y/Z, velX/Y/Z, mass, radius
// 1 int:    state  (to skip INACTIVE ghosts in collision kernel)
// attachLocalX/Y/Z NOT included: halo particles are read-only ghosts
\end{lstlisting}

\subsection{Particle Migration}

Occasionally a particle crosses the subdomain boundary (tethered to the ball
as it moves). Migration must transfer \textbf{all 16 SoA arrays} of
\texttt{ParticleSystem}:

\begin{lstlisting}[numbers=none]
posX, posY, posZ,            // 3 float
velX, velY, velZ,            // 3 float
forceX, forceY, forceZ,      // 3 float
mass, radius, wetness,       // 3 float
state,                       // 1 ShellParticleState (int)
attachLocalX, attachLocalY, attachLocalZ  // 3 float -- CRITICAL
\end{lstlisting}

The \texttt{attachLocal*} arrays hold body-frame tether anchors computed
once at capture via \texttt{quatInvRotateVec}. They are \textbf{not
reconstructible} from world-space state; omitting them from migration would
cause tether forces to point in the wrong direction.

\subsection{Modifications to \texttt{copyParamsToGPU}}

\begin{lstlisting}
// Must be called once per device at startup
for (int g = 0; g < numGPUs; g++) {
    cudaSetDevice(g);
    copyParamsToGPU(params); // writes d_params __constant__ on device g
}
\end{lstlisting}

\subsection{P2P Access Initialization}

\texttt{cudaDeviceEnablePeerAccess} operates from the \emph{current} device
context; both directions must be enabled from the correct device:

\begin{lstlisting}
for (int g = 0; g < numGPUs - 1; g++) {
    cudaSetDevice(g);
    cudaDeviceEnablePeerAccess(g + 1, 0); // g -> g+1
    cudaSetDevice(g + 1);
    cudaDeviceEnablePeerAccess(g, 0);     // g+1 -> g
}
\end{lstlisting}

If P2P is unavailable (PCIe without IOMMU bridge), the fallback is to
transit through a pinned host staging buffer (2$\times$ latency, functionally
identical).

\subsection{Small-N Fused Path}

The single-GPU code switches to \texttt{shellPhysicsFusedKernel} when
$\texttt{shellN} \le 2048$. In multi-GPU, the per-GPU count
$\texttt{shellN}/G$ may fall below this threshold even at moderate total~$N$.
The dispatcher must check the \textbf{local} particle count on each GPU and
select the fused path independently per device.

% ============================================================
\section{CUDA-Aware MPI for Clusters}
\label{sec:mpi}

\subsection{Motivation}

The peer-to-peer mechanism described in Section~\ref{sec:code}
(\texttt{cudaMemcpyPeerAsync}) only works between GPUs connected to the
\emph{same} CUDA driver instance, i.e.\ on a single physical node.
Scaling to a multi-node cluster (16, 32, 64+ GPUs) requires
\textbf{MPI} (Message Passing Interface) to transfer data across the network.

\subsection{Classic MPI vs.\ CUDA-Aware MPI}

With \emph{classic} MPI, transferring a GPU buffer to another node requires
four explicit steps:
\begin{enumerate}
  \item \texttt{cudaMemcpy} D$\to$H: copy halo from GPU to host pinned buffer;
  \item \texttt{MPI\_Isend}: send host buffer over the network;
  \item \texttt{MPI\_Irecv}: receive into host buffer on the remote node;
  \item \texttt{cudaMemcpy} H$\to$D: copy received data from host to GPU.
\end{enumerate}
Steps 1 and 4 add two full PCIe traversals per transfer, doubling effective
latency and consuming host CPU time for staging.

\emph{CUDA-aware MPI} (e.g.\ OpenMPI~$\ge$1.7 with \texttt{--with-cuda}, or
MVAPICH2-GDR) recognises device pointers passed directly to
\texttt{MPI\_Isend}/\texttt{MPI\_Irecv} and, when \textbf{GPUDirect RDMA}
is available (NVLink or RDMA-capable NIC such as Mellanox InfiniBand), it
transfers data directly from GPU to network card without involving the host
CPU or host RAM at all.

\begin{table}[ht]
  \caption{Classic MPI vs.\ CUDA-aware MPI for halo exchange}
  \label{tab:mpi}
  \centering
  \begin{tabular}{lcc}
    \toprule
    & Classic MPI & CUDA-Aware MPI \\
    \midrule
    Data path     & GPU$\to$CPU$\to$NIC$\to$CPU$\to$GPU & GPU$\to$NIC$\to$GPU \\
    Extra copies  & 2 (D$\to$H + H$\to$D) & 0 \\
    Host staging buffer & required & not required \\
    Latency       & $\sim$2$\times$ higher & baseline \\
    Code change   & host copy + MPI call & MPI call only \\
    \bottomrule
  \end{tabular}
\end{table}

\subsection{Implementation}

With CUDA-aware MPI, device pointers are passed directly to the MPI
primitives. The halo pack/unpack kernels are unchanged; only the transfer
call differs:

\begin{lstlisting}[caption={Halo exchange with CUDA-aware MPI (non-blocking)}]
launchPackHalo(gpuData, commStream); // fill haloSendBuf on GPU

// Device pointers passed directly -- no host staging copy
MPI_Isend(gpuData.haloSendBuf, haloBytes, MPI_BYTE,
          neighborRank, tag, MPI_COMM_WORLD, &sendReq);
MPI_Irecv(gpuData.haloRecvBuf, haloBytes, MPI_BYTE,
          neighborRank, tag, MPI_COMM_WORLD, &recvReq);

// Overlap network transfer with forceTetherKernel on simStream,
// then wait before neighborCollisionKernel
MPI_Waitall(2, reqs, statuses);
launchUnpackHalo(gpuData, commStream);
\end{lstlisting}

The captured-mass reduce across all nodes maps to a single MPI collective,
avoiding any manual reduction loop:

\begin{lstlisting}[caption={Captured-mass reduce with MPI}]
// Each rank holds its local accumulated mass after captureKernel
MPI_Allreduce(&localCaptureMass, &totalCaptureMass,
              1, MPI_FLOAT, MPI_SUM, MPI_COMM_WORLD);
// totalCaptureMass now available on all ranks for ball update
updateCoreMassAfterCapture(ball, totalCaptureCount,
                           totalCaptureMass, snowDensity);
\end{lstlisting}

\subsection{Deployment Considerations}

If GPUDirect RDMA is unavailable, CUDA-aware MPI falls back transparently to
the classic D$\to$H$\to$network$\to$H$\to$D path, so the code remains
correct and portable across all cluster configurations.

% ============================================================
\section{Conclusions}
\label{sec:conclusions}

\begin{enumerate}
  \item \textbf{Highly favorable for multi-GPU}: \texttt{neighborCollisionKernel}
        is both the heaviest stage and the only one requiring significant
        communication, giving a comm/compute ratio of $< 0.1\%$.

  \item \textbf{1D domain decomposition along $s$} is $\sim$25$\times$ more
        communication-efficient than index-based partitioning, exploiting the
        strong spatial coherence of shell particles around the ball.

  \item \textbf{Halo width} must be $\max(\texttt{cellSize},
        \texttt{cohesionRadius})$, computed at runtime to remain correct under
        all CLI configurations.

  \item \textbf{Migration must transfer all 16 SoA arrays}, especially
        \texttt{attachLocal*} (body-frame tether anchors, computed once at
        capture and not reconstructible from world-space state).

  \item \textbf{Expected scaling}: nearly linear up to 4--8~GPUs;
        $\sim$3.8$\times$ speedup at 4~GPUs with $N = 200$\,k active shell
        particles.

  \item \textbf{Small-N fused path}: the \texttt{shellPhysicsFusedKernel}
        threshold must be evaluated independently per device in multi-GPU,
        since the local $N/G$ may drop below 2048 even when total $N$ is
        large.

  \item \textbf{Minimal code impact}: changes are confined to \texttt{main.cu}
        (multi-device loop + capture reduce), two new lightweight kernels
        (\texttt{packHalo}, \texttt{unpackHalo}), one additional stream per
        GPU, and grid-sizing adjustments. The existing physics kernels require
        \textbf{no modification}.
\end{enumerate}

\end{document}
